\documentclass{article}
\usepackage[left=3cm, right=3cm, top=1cm, bottom=4cm]{geometry} % Utilizo este paquete para ajustar los márgentes que tenía en word
\usepackage[spanish,es-nodecimaldot]{babel}
\usepackage[utf8]{inputenc} % Este paquete me ayuda a tener soporte para letras y acentos en español
\usepackage{graphicx} % Este me sireve para poder poner imágenes
\usepackage{wrapfig} % Este me sirve para poner las imágenes al costado de los párrafos
\usepackage{pdfpages} % Este me sirve para poder usar pdfs
\usepackage{indentfirst} % Este me permite modificar las sangrías
\usepackage{url} % Este me sirve para poder poner urls
\usepackage{fancyhdr} % Este me permite establecer el estilo de página (osea encabezad y pié de pág)
\usepackage{xcolor}
\usepackage{soul}
\sethlcolor{yellow}
\usepackage{fvextra}
\usepackage{csquotes} % Este me sirve para poner comillas
\usepackage{amsmath} % Este me sirve para poder escribir ecuaciones
\usepackage{amssymb} % Creo que me ayuda con símbolos matemáticos
\usepackage{braket} % Para poder usar la notación de Dirac
\usepackage{enumitem} % Este me sirve para poder dar formato a las listas
\usepackage{cite}
\usepackage{multicol} % Este me sirve para formatear las columnas que quiero en el paper
\usepackage{float}
\usepackage{hyperref}% Este sirve para hacer links
\usepackage[spanish]{cleveref}% Me sirve para los \cref

%Traducciones de cref
\crefname{equation}{ecuación}{ecuaciones}
\Crefname{equation}{Ecuación}{Ecuaciones}
\crefname{figure}{figura}{figuras}
\Crefname{figure}{Figura}{Figuras}
\crefname{table}{tabla}{tablas}
\Crefname{table}{Tabla}{Tablas}
\crefname{section}{sección}{secciones}
\Crefname{section}{Sección}{Secciones}

\title{Ecuación de Schrödinger independiente del tiempo en un potencial central. El Átomo de Hidrógeno}
\author{Lautaro Andres Heredia}
\date{}

\begin{document}
    \maketitle
    
    \newpage
    \tableofcontents
    \newpage
    
    \section{Introducción}
En este trabajo se tratará la Ecuación de Schrödinger independiente del tiempo en un potencial con simetría esférica. En particular se trabajará con el caso del potencial central ( de la forma \(V = V(r)\), donde \(r = |\mathbf{r}|\)) y plantearemos el problema como una ecuación de autovalores. Una vez planteado el problema del potencial central, intentaremos encontrar la solución de la ecuación de Schrödinger independiente del tiempo para el caso del potencial Coulombiano, correspondiente al conocido problema del átomo de Hidrógeno.

%    \input{formalismo}
%    \input{problema}
%    \input{separación de variables/separacion de variables}
%    \input{Átomo de Hidrógeno/Átomo de Hidrógeno}

%    \newpage
%    \nocite{*}
%    \bibliographystyle{plain}
%    \bibliography{referencias}
%    \input{bitacora de ia}
\end{document}
